% Part: counterfactuals
% Chapter: introduction
% Section: counterfactuals

\documentclass[../../../include/open-logic-section]{subfiles}

\begin{document}

\olfileid[hi]{cnt}{int}{cnt}

\olsection{प्रतितथ्यात्मक सशर्त}

अंग्रेज़ी में ``यदि \dots तो \dots'' रचनाओं का एक अत्यंत सामान्य और
महत्त्वपूर्ण रूप \emph{to be} के भूतकालीन संभावनार्थक रूप से बनाया जाता है:
``यदि ऐसा होता कि \dots तो ऐसा होता कि \dots।'' चूँकि ऐसे सशर्त का
पूर्ववर्ती सामान्यतः असत्य, अर्थात तथ्य के विपरीत, होता है, इन्हें
\emph{प्रतितथ्यात्मक सशर्त} कहते हैं (और \emph{to be} के संभावनार्थक रूप के
प्रयोग के कारण \emph{संभावनार्थक सशर्त} भी)। इन्हें उन \emph{विधानवाचक}
सशर्तों से अलग किया जाता है जिनका रूप ``यदि ऐसा है कि \dots तो ऐसा है कि
\dots'' होता है। प्रतितथ्यात्मक और विधानवाचक सशर्तों की सत्यता-शर्तें भिन्न
हैं। एडम्स के प्रसिद्ध उदाहरण पर विचार करें:
\begin{quote}
  यदि ऑज़वाल्ड ने केनेडी को नहीं मारा, तो किसी और ने मारा।
  
  यदि ऑज़वाल्ड ने केनेडी को न मारा होता, तो किसी और ने मारा होता।
\end{quote}
पहला विधानवाचक और दूसरा प्रतितथ्यात्मक है। पहला स्पष्टतः सत्य है: हम जानते
हैं कि राष्ट्रपति जॉन एफ. केनेडी को \emph{किसी} ने मारा था; और यदि वह
व्यक्ति (वॉरेन रिपोर्ट के विपरीत) ली हार्वी ऑज़वाल्ड नहीं था, तो किसी अन्य
ने केनेडी को मारा। दूसरा कुछ अलग कहता है। वह दावा करता है कि यदि ऑज़वाल्ड
ने केनेडी को न मारा होता---अर्थात डैलस की गोलीबारी टल गई होती या असफल रही
होती---तो बाद का इतिहास इस प्रकार घटित होता कि हत्या का कोई दूसरा प्रयास
सफल हो जाता। उसके सत्य होने के लिए शक्तिशाली तत्वों ने JFK की मृत्यु
सुनिश्चित करने का षड्यंत्र किया होना चाहिए (जैसा अनेक JFK षड्यंत्र-सिद्धांतकार
मानते हैं)।

इस पर सक्रिय बहस है कि क्या \emph{विधानवाचक} सशर्त को भौतिक सशर्त ठीक से
ग्रहण करता है; विशेषतः क्या भौतिक सशर्त के विरोधाभासों को इस प्रकार
``समझाया'' जा सकता है जो अंग्रेज़ी विधानवाचक सशर्तों की सत्यता-शर्तें देने
के उसके दावे से संगत हो। इसके विपरीत यह निर्विवाद है कि प्रतितथ्यात्मक सशर्तों
को भौतिक सशर्त द्वारा ठीक से संकेतित नहीं किया जा सकता। कारण स्पष्ट है:
यद्यपि प्रतितथ्यात्मक सशर्तों के पूर्ववर्ती सामान्यतः असत्य होते हैं, असत्य
पूर्ववर्ती वाला हर प्रतितथ्यात्मक सशर्त सत्य नहीं होता। उदाहरणार्थ, यदि आप
वॉरेन रिपोर्ट मानते हैं और JFK की हत्या का कोई षड्यंत्र नहीं था, तो एडम्स का
प्रतितथ्यात्मक सशर्त इसका उदाहरण है।

कारणात्मक तर्क-विचार में प्रतितथ्यात्मक सशर्त महत्त्वपूर्ण भूमिका निभाते हैं:
प्रतितथ्यात्मक के उपयोग का प्रमुख उदाहरण कारणात्मक संबंध व्यक्त करना है।
मसलन, माचिस की तीली रगड़ने से वह जलती है; इसे हम कह सकते हैं, ``यदि इस तीली
को रगड़ा जाता, तो वह जलती।'' भौतिक सशर्त और सामान्यतः विधानवाचक सशर्त इसे
व्यक्त नहीं कर सकते: ``तीली रगड़ी जाती है $\lif$ तीली जलती है'' तब सत्य है
जब तीली कभी रगड़ी ही न जाए, चाहे रगड़े जाने पर कुछ भी होता। इससे भी बुरा,
``तीली रगड़ी जाती है $\lif$ तीली फूलों का गुलदस्ता बन जाती है'' भी तब सत्य
है जब उसे कभी न रगड़ा जाए; किंतु रगड़े जाने पर तीली निश्चित रूप से फूलों का
गुलदस्ता नहीं बनेगी।

प्रतितथ्यात्मक का सही तर्कशास्त्र ठीक क्या है, इस पर अभी बहस है। स्टालनेकर
और लुईस ने प्रतितथ्यात्मक का प्रभावशाली विश्लेषण दिया। उनके अनुसार कोई
प्रतितथ्यात्मक ``यदि~$S$ होता, तो~$T$ होता'' तभी सत्य है जब~$T$ उस
प्रतितथ्यात्मक स्थिति (``संभव जगत'') में सत्य हो जो वास्तविक जगत की स्थिति
के सबसे निकट है और जहाँ~$S$ सत्य है। इसे ``सत्तामीमांसात्मक'' विश्लेषण कहते
हैं, क्योंकि यह संभव जगतों की सत्तामीमांसा का उल्लेख करता है। अन्य विश्लेषण
सशर्त प्रायिकताओं या विश्वास-संशोधन के सिद्धांतों का उपयोग करते हैं।
प्रतितथ्यात्मक के अनेक भिन्न प्रस्तावित तर्कशास्त्र हैं। प्रतितथ्यात्मक का
कोई एक लुईस--स्टालनेकर तर्कशास्त्र तक नहीं है: यद्यपि स्टालनेकर और लुईस ने
निकटतम संभव जगतों के संदर्भ में समान दिशा के विवरण प्रस्तावित किए, उनकी
मान्यताएँ अलग-अलग मान्य अनुमान देती हैं।

\end{document}
